\section{Changing the fiber center}\label{sec:centers}

A moving center is not a holomorphic change of coordinates. We therefore first widen the fibers, then estimate the actual change of potential as a perturbation in fixed holomorphic coordinates. The widening and the metric scales must be chosen in the order specified below.

Throughout this section $p=\kappa=1$ and $Y\ge\max\{1,t\}$, so $\lambda,H$ are constant by Lemma~\ref{radial:odes}, and $V=\lambda G_0$ as metrics, including their jets. We use the scaled affine spatial coordinates of Lemma~\ref{radial:scaled}. The severity $E$ is constant, and no shear term is present. The old twist has the form
\[
 \psi_{\rm old}=\Psi+\psi_e,
 \qquad G_e=\partial\bar\partial\psi_e,
 \qquad |\partial^{[j]}G_e|\le\delta/H\quad(0\le j\le2),
\]
with $\delta$ sufficiently small as in Proposition~\ref{radial:aligned}. Suppose also that, throughout the interval being configured,
\begin{equation}\label{center:old-data}
 \psi_{\rm old}\le A_{\rm old},\qquad
 |\partial^{[j]}\psi_{\rm old}|\le A_{\rm old}
 \quad(1\le j\le4)
\end{equation}
in these fixed scaled axes at each point. The number $A_{\rm old}$ is known independently of the target sizes of $\lambda,H$. Additive constants already placed in the twist are included in this upper bound.

\begin{proposition}[Center change]\label{center:change}
Choose the profile constants in the following order: first the fixed exponent $b_0$ of the quiet profile; then a sufficiently large absolute constant $d_0$; and then $K_0$ sufficiently large in terms of $d_0$ and absolute constants specified in the proof. Fix $E$ and old data as in \eqref{center:old-data}. There is a width parameter $L$, chosen using only these data and the fixed profile, such that the following construction is possible.

After choosing $L$, choose a sufficiently large $\lambda$, and then a sufficiently large $H$; explicitly one can require
\begin{equation}\label{center:target-order}
 \lambda\ge\max\{e^{3L},\Lambda_f\},
 \qquad H\ge M_f(1+\log\lambda)^2,
\end{equation}
where $\Lambda_f,M_f$ are finite constants determined by the fixed profile, $E,L$, the old data, and $\delta$, but not by $\lambda,H$. Over three fixed units of the variable $\Psi$, one can change from any prescribed constant center in the unit disk to any other such center, returning to the original profile $q_0(x,E)$ at the end. The final twist differs from the old twist by a nonnegative constant, possibly very large.

During the change the potential is smooth, its metric is positive and has nonpositive real sectional curvature, and the metric bounds $cV$ from below on spatial projection, for a fixed $c>0$. There is no restriction on the cost of the final constant offset. All power exponents and smallness tolerances used before choosing the spatial dimension are dimension independent.
\end{proposition}

\subsection{The profile deformation and its elementary bounds}

Write $F^0_x=1/q_0$ and normalize $F^0(x)-x\to0$ at the left end. We record properties of the concrete quiet profile that will be used here. At fixed $E$,
\begin{equation}\label{center:q-properties}
 \begin{split}
 &q_0\ge1,\quad (q_0)_x>0,\quad
       (q_0)_{xx}/(q_0)_x\ge-1/16,\\
 &|(q_0)_x|+|(q_0)_{xx}|\le C q_0,\qquad
       q_0(x,E)\le C_E+C(1+\max(x,0)),\\
 &\inf_{x\ge-20}(q_0)_x>0.
 \end{split}
\end{equation}
Here and below a subscript $E$ permits dependence on $E$ and the already fixed profile. The last infimum can be bounded below using just the derivative of $K_0^{-1}\log(1+e^x)$. The inequality for the second derivative follows by differentiating both summands in the quiet profile. In particular $(q_0)_x+(q_0)_{xx}\ge(15/16)(q_0)_x$.

Two endpoint properties are also needed. The function $Q_0(v)=q_0(\log v,E)$ is smooth down to $v=0$, with
\begin{equation}\label{center:left-slope}
 Q_0(0)=1,\qquad Q_0'(0)=b_0e^E+K_0^{-1}>0.
\end{equation}
At the other end the auxiliary function in the quiet profile tends to a positive finite limit with error $O(e^{-x/16})$, so
\[
 q_0(x,E)=c_E+x/K_0+O_E(e^{-x/16})\quad(x\longrightarrow\infty).
\]
Integration, using $1/q_0-K_0/x=O_E(x^{-2})$, gives
\begin{equation}\label{center:primitive-asymptotic}
 F^0(x)=K_0\log x+O_E(1),\qquad
 \int_0^L\frac{dx}{q_0(x,E)}=K_0\log L+O_E(1).
\end{equation}
No uniformity in $E$ is asserted or needed in these endpoint constants; $E$ has been fixed before $L$.

Take $L$ large. Choose a smooth function $\chi$ with
\[
 0\le\chi\le1,\quad \operatorname{supp}\chi\subset[-20,3L],\quad
 \chi=1\text{ on }[0,L],\quad \int\chi=2L,\quad \chi'\ge-1/16,
\]
and with its derivatives through any fixed needed order bounded absolutely. Such a function is obtained from a fixed rising cutoff, an adjustable plateau, and a fixed descending cutoff of length greater than $16$; adjusting the plateau gives the integral exactly. Define the increasing inverse map $m=m(x)$ by
\begin{equation}\label{center:inverse-map}
 x=m+L-\frac12\int_{-\infty}^m\chi(v)\,dv.
\end{equation}
Then
\begin{equation}\label{center:map-bounds}
 1\le m'\le2,\quad m\le x,\quad
 m=x-L\ (x\le L-20),\quad m=x\ (x\ge X:=3L).
\end{equation}
Its derivatives of any fixed order are bounded absolutely. Set
\[
 q_1(x)=q_0(m(x),E)\le q_0(x,E).
\]
Since $m''=\chi'(m)(m')^3/2$, putting $a=m'\in[1,2]$ gives
\begin{align}\label{center:q1-positive}
 (q_1)_x+(q_1)_{xx}
 &\ge(q_0)_x(m)\bigl(a-a^2/16-a^3/32\bigr)\\
 &\ge\tfrac34(q_0)_x(m)a=\tfrac34(q_1)_x>0.\notag
\end{align}
Also $|(q_1)_x|+|(q_1)_{xx}|\le Cq_1$, with constants independent of $L$. These bounds follow from \eqref{center:q-properties} and the bounds on $m',m''$.

\subsection{Offsets during profile interpolation}

On the first unit of $\Psi$, choose a smooth nondecreasing flat-ended cutoff $\beta=\beta(\Psi)$ from $0$ to $1$, with fixed bounds on its derivatives through order four. Define
\begin{equation}\label{center:forward-profile}
 q(x,\beta)=(1-\beta)q_0(x,E)+\beta q_1(x),\qquad
 \mathfrak s=1/q,
\end{equation}
and let $F^\beta$ have derivative $F^\beta_x=\mathfrak s$ and left normalization $F^\beta-x\to0$. Every such primitive exists, tends to infinity on the right, and differs from $F^0$ by a constant for $x\ge X$. All spatial derivatives in the rest of this subsection are at fixed $\tau=e^x$.

Use an offset $w^+(\beta)$ defined by $w^+(0)=0$ and
\begin{equation}\label{center:offset-definition}
 (w^+)'(\beta)=\partial_\beta F^\beta(X)
                  +\frac{d_0}{1-\beta+1/L},\qquad
 D(x,\beta)=(w^+)'(\beta)-\partial_\beta F^\beta(x).
\end{equation}
Here the prime on $w^+$ denotes differentiation in $\beta$. Since $\mathfrak s_\beta\ge0$ and vanishes for $x\ge X$,
\begin{equation}\label{center:D-domination}
 D\ge\frac{d_0}{1-\beta+1/L}.
\end{equation}
This lower bound can dominate $\sup_x|\mathfrak s_\beta|$ by any required fixed factor if $d_0$ is sufficiently large. Indeed
\[
 0\le\mathfrak s_\beta=(q_0-q_1)/q^2
 \le\min\{(1-\beta)^{-1},\sup_{x\le X}q_0(x,E)\}.
\]
By taking $L$ sufficiently large relative to $C_E$, the second entry is at most $CL$ with an absolute $C$. The elementary bound
$\min\{z^{-1},CL\}\le(C+1)/(z+1/L)$ proves the assertion, including $\beta=1$ by continuity.

For later restoration use, on another unit of $\Psi$, the reverse mixture $q=(1-\beta)q_1+\beta q_0$ with a fresh increasing cutoff. Now $\mathfrak s_\beta\le0$ and $\partial_\beta F^\beta\le0$. Choose a new offset $w^-(\beta)$ with $w^-(0)=0$ and constant positive derivative greater than a sufficiently large fixed multiple of $1+\sup_{x,\beta}|\mathfrak s_\beta|$. With $D=(w^-)'-\partial_\beta F^\beta$, the same domination holds. No estimate on the size of $w^-(1)$ is required.

Once $L$ is fixed, all the following constants are finite and may be included in a number $C_f$: upper bounds for the offsets and their derivatives through order four; for $D$ and its first three $\beta$ derivatives; for the profile and the finite-order derivatives used below on $x\le X$; for their left-end decay coefficients; and the reciprocal lower constants in the positive estimates on that interval. To check finiteness at the noncompact left endpoint, both $q_0-1$ and $q_1-1$, with their needed derivatives, are $O_f(e^x)$. Thus $\beta$ derivatives of $\mathfrak s$ are $O_f(e^x)$, and their integrals defining derivatives of $F^\beta$ converge. In particular $C_f$ may include $e^X$ and large functions of $E,L$. It is selected before $\lambda,H$.

\begin{lemma}[Curvature during interpolation]\label{center:interpolation}
During either profile switch above, keep the center constant and set
$\psi=\psi_{\rm old}+w^\pm(\beta(\Psi))$, also including the already frozen offset when making the reverse switch. If $\lambda$ and then $H$ satisfy sufficiently large thresholds of the form \eqref{center:target-order}, the metric is positive, $h\ge cV$, and its negative curvature has a strictly positive rank-two reserve off the fiber axis. These conclusions extend across that axis by continuity.
\end{lemma}

\begin{proof}
A prime on $\beta$ now means differentiation in $\Psi$. Direct differentiation of $\psi-F^\beta$ gives
\begin{equation}\label{center:N-formula}
 \mathcal N=(1+D\beta')G_0+G_e
       +\bigl(D\beta''+D_\beta(\beta')^2\bigr)
                   \partial\Psi\otimes\bar\partial\Psi.
\end{equation}
Here $D_\beta$ is the derivative at fixed $x$; thus all dependence of the primitive on the spatial cutoff is included. The last term and its first two spatial jets are $O(C_f/H^2)$, because all positive derivatives of $\Psi$ in scaled axes are $O(1/H)$. The derivatives of the scalar $D\beta'$ beyond its value are $O(C_f/H)$. With $d_f=1+D\beta'\ge1$, these facts show
\[
 \mathcal N\asymp d_fG_0>0,
 \qquad h=V+\int_0^\tau\mathcal N(v)\,dv\ge cV.
\]
The constants in $\asymp$ can be absolute by making $H$ large enough; the value $d_f$ itself is bounded by $C_f$.

First consider $x\le X$. Since $\tau\le e^X$, the first two metric jets of $h-V$ are $O(C_f/H)$. Hence
\begin{equation}\label{center:finite-x-h}
 \Gamma^h=\Gamma^{G_0}+O(C_f/\lambda),\qquad
 C^h(T,T)\ge c\lambda|T|^2/H
 \quad(\operatorname{rank}T\le2).
\end{equation}
The second estimate follows from Lemma~\ref{radial:curvature}, the scaled jet estimates, and the curvature perturbation identity, taking $\lambda\gg C_f$.

The mixed form in Lemma~\ref{fiber:blocks} is
\[
 \frac{W}{Pq}=(1+r)\mathcal N
       -\tau\mathcal N h^{-1}\mathcal N
       +q^{-1}\partial\bar\partial\log q.
\]
The last term equals
\[
 -\mathfrak s_\beta\beta'G_0+O(C_f/H^2).
\]
The spin form, divided by $Pq$, equals
\[
 -\mathfrak s_\beta\beta'\nabla_h^{2,0}\Psi+O(C_f/H^2).
\]
In fact $q_\beta/q^2=-\mathfrak s_\beta$; the remaining terms contain two first derivatives of $\Psi$. Lemma~\ref{radial:scaled} and \eqref{center:finite-x-h} give
$|\nabla_h^{2,0}\Psi|\le C/H$ with absolute $C$. Also
$|\tau\mathcal N h^{-1}\mathcal N|\le C_f/(\lambda H)$. Choose $d_0$ first so that the domination in \eqref{center:D-domination} is sufficiently strong in comparison with this absolute $C$, and then choose $\lambda,H$ large compared with $C_f$. It follows that
\begin{equation}\label{center:W-spin}
 W\ge c\frac{Pq}{H}d_f I,
 \qquad |\nabla_h^{2,0}P(b,b)|\le\tfrac14 W(b,\bar b).
\end{equation}
This argument also applies to the reverse mixture, using absolute values of $\mathfrak s_\beta$.

By \eqref{center:q1-positive} and convexity of the mixture,
\[
 D_q=(q_x+q_{xx})/q\ge c_f\min(1,\tau)\quad(x\le X).
\]
Left-end decay and boundedness on the rest of this interval give
\[
 |\partial\log P|+|\partial P_\tau|
       \le C_f\min(1,\tau)/H.
\]
It follows from the curvature table that
\begin{equation}\label{center:A-control}
 A=(1+o(1))Pq^2D_q,\qquad
 |A_i|\le C_f P\min(1,\tau)/H,
\end{equation}
where the relative error tends uniformly to zero as $\lambda,H$ increase. To verify uniformity at $\tau=0$, the subtracted term in $A$ is at most
$C_fP^2\min(1,\tau)^2/(\lambda H)$, while the leading term is $Pq^2D_q$. The ratio tends to zero uniformly under the stated choices. Similarly \eqref{center:W-spin} and \eqref{center:A-control} make the $A_i$ coupling arbitrarily small in units $(AW)^{1/2}$.

Finally $|K/P|\le C_f/H$. Spending fixed portions of $A,W$ according to Lemma~\ref{fiber:ranktwo} costs at most
\[
 C_f\frac{P}{qH}|T|^2\le\frac{C_f}{H}|T|^2
 \quad\text{for the }K\text{ coupling},\qquad
 C_f\frac{P}{D_qH^2}|T|^2\le\frac{C_f}{H^2}|T|^2
 \quad\text{for the trace coupling}.
\]
In these displayed bounds $C_f$ is enlarged to include $e^X$; also $P/D_q$ is bounded on $x\le X$, including the left endpoint. Both costs are small compared with \eqref{center:finite-x-h}. Thus fixed fractions of $A$, $W$, and $\lambda|T|^2/H$ remain.

Now let $x>X$. Since the two profiles coincide exactly here, $F^\beta(x)-F^0(x)$ is independent of $x$. Consequently $\mathcal N=\mathcal N_\infty$ is independent of $\tau$, and
\begin{equation}\label{center:finite-tail}
 h=V_f+\tau\mathcal N_\infty,
 \qquad V_f=V+\int_0^{e^X}
          \bigl(\mathcal N(v)-\mathcal N_\infty\bigr)\,dv.
\end{equation}
There is no infinite tail in this expression. Its first two spatial jets differ from those of $V$ by $O(C_f/H)$. Treating $d_f=1+D\beta'$ at the point as constant, the jets of $\mathcal N_\infty$ differ from those of $d_fG_0$ by at most $C\delta d_f/H$, after taking $H\gg C_f/\delta$; similarly those of $V_f$ differ from $\lambda G_0$ by a sufficiently small relative amount after taking $\lambda\gg C_f/\delta$.

The aligned argument can now be applied with the two spatial scales $\lambda$ and $d_f$. More explicitly, let $\theta_f=\lambda/(\lambda+\tau d_f)$. The comparisons and the identity $\nabla^h(V_f+\tau\mathcal N_\infty)=0$ give
\[
 C^h(T,T)\ge c\frac{\lambda+\tau d_f}{H}|T|^2,
 \qquad |\nabla^h\mathcal N_\infty|
                   \le C\delta d_f\theta_f/H.
\]
The $K$ cost is bounded by $C\delta^2(\lambda+\tau d_f)|T|^2/H$. The $rI$ part of the trace cost uses the same rank-two symmetrized $P\mathcal N_\infty^2$ reserve as before. The remaining cost is at most
\[
 C\frac{\tau qd_f^2\theta_f^2}{D_qH^2}|T|^2
 \le C_f\frac{(1+\log\lambda)^2(\lambda+\tau d_f)}{H^2}|T|^2.
\]
For the last step, use \eqref{radial:profile-bounds} and the scalar inequality from the proof of Proposition~\ref{radial:aligned} with $\lambda/d_f$ in place of $\lambda$. We have $1\le d_f\le C_f$ and may require $\lambda\ge C_f$, so the logarithm is bounded by $\log\lambda$. Taking $H\ge M_f(1+\log\lambda)^2$ with sufficiently large $M_f$ absorbs this last cost. Spin and $A_i$ vanish in this region because $q=q_0$ is spatially constant.

We have proved positivity and nonpositive curvature off the axis, with fixed fractional reserves. The profiles are smooth in $\tau$ at zero, equal to one there, and have the prescribed left normalization. Thus the potential and metric are smooth in the original affine fiber coordinate, whose metric at the axis has entries $V,e^\psi$. Curvature extends by continuity. This proves the lemma.
\end{proof}

\subsection{The width gained before moving the center}

At the end of the first switch freeze $\beta=1$ and the offset, and write
$\psi=\psi_{\rm old}+w^+(1)$. Changing variables in the defining integral gives
\begin{equation}\label{center:F1-formula}
 F^1(x)=L+F^0(m(x))
          -\frac12\int_{-\infty}^{m(x)}\frac{\chi(v)}{q_0(v,E)}\,dv.
\end{equation}
Both sides have derivative $1/q_1$ and the same left normalization. At $x=X$, $m(X)=X$, so integration of \eqref{center:offset-definition} yields
\begin{equation}\label{center:offset-total}
 w^+(1)=L-\frac12\int_{\mathbb R}\frac{\chi(v)}{q_0(v,E)}\,dv
                  +d_0\log(1+L).
\end{equation}
In particular this offset is nonnegative: the integral is at most $\int\chi=2L$.

Define the pointwise width scales
\[
 R_0=\exp\bigl((L-\psi)/2\bigr),\qquad
 Z=\exp\bigl((F^1(x)-\psi)/2\bigr)=|y-\zeta|.
\]
Because $\chi=1$ on $[0,L]$, \eqref{center:primitive-asymptotic} and \eqref{center:offset-total} imply
\begin{equation}\label{center:R-growth}
 R_0\ge c_{E,A_{\rm old}}L^{K_0/4-d_0/2}.
\end{equation}
Choose $K_0$ so large relative to $d_0$ that $K_0/4-d_0/2>K_0/9$. Then, for $L$ sufficiently large depending on old data,
\begin{equation}\label{center:width-bounds}
 R_0\ge L^{K_0/9},\qquad
 Z\ge x^{K_0/9}\quad(x\ge L-20).
\end{equation}
Here is the second bound in detail. At $x=L-20$, formula \eqref{center:F1-formula} gives $F^1=L+F^0(-20)$, so $Z$ is a fixed old-data multiple of $R_0$. Monotonicity of $F^1$ and \eqref{center:R-growth} prove the desired bound on $[L-20,3L]$, absorbing the fixed factor $3^{K_0/9}$ by increasing $L$. For $x\ge3L$, the integral in \eqref{center:F1-formula} has reached its total and cancels the integral in \eqref{center:offset-total}. Thus exactly
\[
 F^1(x)-\psi=F^0(x)-\psi_{\rm old}-d_0\log(1+L).
\]
The asymptotic \eqref{center:primitive-asymptotic}, together with $x\ge3L$ and the strict exponent inequality already chosen, gives the second bound on this remaining interval. These estimates are uniform over the spatial interval because only the upper bound in \eqref{center:old-data} was used.

\subsection{The actual nonholomorphic center motion}

The preceding calculation has produced width before any center is moved.
We now spend that width to control the metric and curvature errors of
moving the center. Two coordinate systems cover the fiber: logarithmic
coordinates away from the enlarged core and a scaled affine coordinate
on the core, including its axis.

On the next unit of $\Psi$, hold $q=q_1$ and all offsets fixed, and let $\zeta$ be a smooth flat-ended convex-combination interpolation of the two prescribed centers. Its derivatives through order four as a function of $\Psi$ are bounded absolutely. By Lemma~\ref{radial:scaled}, and $H\ge1$, its spatial derivatives in scaled axes satisfy
\begin{equation}\label{center:zeta-jets}
 |\partial^{[j]}\zeta|\le C/H\qquad(1\le j\le4).
\end{equation}
At a point $s_0$, compare the moving-center potential with the potential obtained by replacing $\zeta(s)$ by the constant $\zeta(s_0)$ in a neighborhood. This is a comparison in fixed holomorphic coordinates, not a claim that $y-\zeta(s)$ is a holomorphic coordinate. All constants in the following two comparisons may depend on old data, $E$, and the fixed profile, but are independent of $L,\lambda,H$ once the stated thresholds hold.

First suppose $x\ge L-20$. For the frozen center use scaled spatial axes and the tilted logarithmic fiber coordinate of Lemma~\ref{fiber:blocks}. The previous interpolation proof with cutoff constant retains fractions of $A,W$ and the horizontal curvature reserve. On this region $m\ge-20$, so \eqref{center:q-properties} and \eqref{center:q1-positive} give
\[
 A\ge c\tau,\qquad W\ge c\tau H^{-1}I.
\]
The horizontal reserve is at least $c\tau/H$: on $x\le X$ use $\lambda\ge e^X\ge\tau$ in \eqref{center:finite-x-h}, and on $x>X$ use the aligned reserve. Consequently the full rank-two curvature contraction is bounded below by $c\tau/H$ times the squared Frobenius norm of the test. The metric has the same lower scale and its inverse has norm at most $CH/\tau$. These assertions also follow directly for the metric from $h=V+\tau G$ and $P=\tau q$.

The polynomial holomorphic tilt has coefficients bounded by the first two old twist derivatives; the offset is constant. In these tilted coordinates the frozen connection has norm at most $Cq$. Indeed the possibly large entries in the cubic table are $P_\tau=q+q_x=O(q)$ and $P\mathcal Nh^{-1}=O(q)$, whereas $P_i=0$, the tilted pure holomorphic second spatial derivative vanishes, and $\Gamma^h=O(1)$.

Let $\Phi$ be the one-variable fiber potential as a function of
$\rho+\psi=\log|y-\zeta|^2+\psi$. Its first derivative is $\tau$, and differentiation in this argument is $q\partial_x$. Thus its derivatives of orders one through four have size at most $C\tau q^3$. For example
\[
 \Phi''=\tau q,\quad
 \Phi'''=\tau q(q+q_x),\quad
 \Phi''''=\tau q\bigl((q+q_x)^2+q(q_x+q_{xx})\bigr).
\]
The fixed argument has bounded jets through order four in the chosen frame. The difference in its argument after moving the center is
\begin{equation}\label{center:log-argument-difference}
 \log\left|1-\frac{\zeta(s)-\zeta(s_0)}{y-\zeta(s_0)}\right|^2.
\end{equation}
It vanishes at the evaluation point, and each positive jet contains a derivative of $\zeta$ and a denominator of size at least $Z$. By \eqref{center:zeta-jets}, bounded tilt coefficients, and $HZ\ge1$, its jets through order four are $O(1/(HZ))$. The chain rule consequently bounds every potential-jet difference through order four by
\begin{equation}\label{center:outer-jet-error}
 C\frac{\tau q^3}{HZ}.
\end{equation}
In particular its relative metric error is at most $Cq^3/Z$. Once this is small, Lemma~\ref{fiber:perturbation}, the bound $\Gamma=O(q)$, and the inverse bound $O(H/\tau)$ show that the curvature error on unit rank-two tests is at most
\begin{equation}\label{center:outer-curvature-error}
 C\frac{\tau}{H}\left(\frac{q^5}{Z}
                            +\frac{q^8}{Z^2}\right)
 \le C\frac{\tau q^8}{HZ}.
\end{equation}
The first term accounts for all linear fourth-jet and cubic-contraction changes, and the second is the quadratic inverse-contraction term. Fixed-rank contraction costs only an absolute factor.

For $x\ge L-20$, \eqref{center:q-properties} and $m\le x$ give $q=q_1\le Cx$ once $L$ is large relative to $C_E$. This constant is independent of $L,\lambda,H$. Choose, in addition to the earlier requirement, $K_0/9>10$. Then \eqref{center:width-bounds} makes the ratios in \eqref{center:outer-jet-error} and \eqref{center:outer-curvature-error} arbitrarily small compared with the metric and curvature reserves by increasing $L$. This choice precedes $\lambda,H$.

It remains to include $x<L-20$ and the axis itself. Put $v=\tau e^{-L}$. By \eqref{center:map-bounds}, the fiber potential added to $B$ is exactly $e^L$ times the original smooth radial function, now evaluated at
\[
 |y-\zeta(s)|^2e^{\psi(s)-L}.
\]
This identity follows by changing variables in $\int_0^\tau q_1(\log t)^{-1}dt$. Use scaled spatial axes and the fixed affine fiber coordinate
\[
 z=\frac{y-\zeta(s_0)}{R_0(s_0)}.
\]
At the evaluation point $|z|^2=e^{F^0(\log v)}\asymp_E v$ for $0\le v\le e^{-20}$. The comparison is uniform down to zero because $F^0(\log v)-\log v$ is smooth there. The smooth radial function and its derivatives on this fixed compact argument interval have bounds depending only on old data and $E$.

For the frozen center, \eqref{center:left-slope} gives, with positive old-data constants,
\[
 A\ge c_E\tau^2e^{-L},\qquad
 W\ge c_E\tau H^{-1}I,
\]
throughout $0<v\le e^{-20}$. For the first inequality, the continuous positive ratio of $P^2P_{\tau\tau}$ to $\tau^2e^{-L}$ has a positive limit at zero by $Q_0'(0)>0$; positivity away from zero follows from \eqref{center:q-properties}. For $W$, the term $Pq\mathcal YG$ has a positive lower bound of this size because $\tau\le e^{L-20}$ and $\lambda\ge e^{3L}$. Divide the vertical and cross curvature blocks by $|z|^4$ and $|z|^2$, respectively, on changing from logarithmic to affine fiber coordinates. The retained vertical and cross reserves are then at least $c_Ee^L$ and $c_Ee^L/H$. The horizontal reserve $c\lambda/H$ also dominates $ce^L/H$. The metric has the same smallest scale, so its inverse is bounded by $C_EH/e^L$.

Undoing the exponentiated tilt changes the affine fiber variable by $z\mapsto e^{\theta(s)}z$, where $\theta(s_0)=0$ and its coefficients are bounded by old twist derivatives. Since $|z|$ is bounded, the differential and its inverse, and the needed higher jets, have old-data bounds independent of $L,\lambda,H$. This preserves the indicated rank-two lower margin up to old-data constants. The connections in these coordinates are bounded: the logarithmic entry transforms to
\[
 (P_\tau-1)/z=O_E(|z|),
\]
and the spatial-fiber entry is bounded by
\[
 |P\mathcal Nh^{-1}|/|z|
       \le C_E\tau/(\lambda|z|)\le C_E|z|,
\]
using $\lambda\ge e^L$. The pure vertical-to-spatial entry vanishes before the bounded tilt is undone. The remaining spatial entries are bounded by the radial jet estimates.

In this affine frame the moving smooth argument, before applying the fixed smooth radial function, is
\begin{equation}\label{center:affine-argument}
 \left|z-\frac{\zeta(s)-\zeta(s_0)}{R_0(s_0)}\right|^2
                         e^{\psi(s)-\psi(s_0)}.
\end{equation}
Its difference from the frozen argument has jets through order four bounded by $C/(HR_0)$, including at $z=0$, by \eqref{center:zeta-jets}. Therefore the potential jets change by at most $Ce^L/(HR_0)$. The inverse metric and bounded connection estimates give a curvature error at most
\begin{equation}\label{center:inner-error}
 C\frac{e^L}{H}\left(R_0^{-1}+R_0^{-2}\right).
\end{equation}
By \eqref{center:width-bounds}, this is a small fraction of the $e^L/H$ reserve for $L$ sufficiently large, independently of the later choices of $\lambda,H$. The metric error is small by the same estimate. Smoothness includes the axis, so no limiting logarithmic-coordinate assertion is being used at a singular coordinate point.

We have proved that the moving-center metric remains, say, at least half the frozen-center metric and that its real sectional curvature is nonpositive. Its spatial projection bound is therefore inherited from the frozen Schur bound $h\ge cV$.

\begin{proof}[Completion of Proposition~\ref{center:change}]
Choose $d_0$ to dominate the absolute spin constants in the interpolation estimate, and choose $K_0$ so that $K_0/4-d_0/2>K_0/9$ and $K_0/9>10$. The latter explicit exponent is larger than the exponent $8$ in \eqref{center:outer-curvature-error}. For fixed old data and $E$, choose $L$ sufficiently large for \eqref{center:width-bounds} and for the moving-center metric and curvature errors to be less than the retained reserves. Next compute the finite constants $C_f$ for both switches and choose $\lambda$ and then $H$ according to \eqref{center:target-order}, with the larger thresholds needed in Lemma~\ref{center:interpolation}. This is the required noncircular order of choices.

Perform the forward switch with constant old center, the center interpolation with frozen profile $q_1$ and offset, and finally the reverse switch with constant new center. Lemma~\ref{center:interpolation} controls the two switches and the preceding argument controls the actual center motion. The cutoffs are flat-ended, so all joins are smooth. At the end the profile is $q_0$, the center is the prescribed new constant, and the added offset is $w^+(1)+w^-(1)\ge0$. All profiles are at least one, have smooth left endpoints with value one, and have normalized primitives tending to infinity on the right; thus they define smooth fiber potentials for all affine fiber values throughout. The spatial projection bound follows on each part, with a fixed fractional loss during the center motion. No bound on the final constant offset is required or asserted.
\end{proof}
